\begin{summarybox}
	\textbf{\textcolor{CSummary}{一句话核心}}
	射影面积法将二面角的余弦绝对值转化为原图形与其正投影的面积比，锐角取正、钝角取负。

	\medskip
	\textbf{\textcolor{CSummary}{使用条件}}
	\begin{itemize}[leftmargin=*,itemsep=0.5em,topsep=0.4em]
		\item\textbf{条件1（图形无重叠）：}原平面图形为简单图形，无折叠、无重叠。\item\textbf{条件2（正投影前提）：}投影必须是正投影（投影线垂直于投影面）。\item\textbf{条件3（角度范围）：}二面角$\theta$满足$0<\theta<\pi$。
	\end{itemize}

	\medskip
	\textbf{\textcolor{CSummary}{关键提醒}}
	\begin{itemize}[leftmargin=*,itemsep=0.5em,topsep=0.4em]
		\item\textbf{易错点（符号判别）：}求出$|\cos\theta|$后，需根据二面角是锐角还是钝角决定$\cos\theta$的正负。\item\textbf{检查项（投影真实性）：}确认投影是否为原图形在另一平面上的正投影，避免误用斜投影。
	\end{itemize}

\end{summarybox}
